\documentclass[11pt]{article}
\usepackage[margin=1in]{geometry}
\usepackage{booktabs}
\usepackage{array}
\usepackage{hyperref}

\title{BSFS Computer Vision Model Progress and Research Direction}
\author{GuTelligence Model Development}
\date{June 5, 2026}

\begin{document}
\maketitle

\section{Executive Summary}

The current best model is a ConvNeXt-Tiny classifier trained on the leakage-safe clean split. Under strict 7-class BSFS classification, performance remains limited. However, when evaluated under a product-aligned grouping strategy, performance improves substantially.

The strongest current product direction is to move from exact 7-class output to a 3-class primary product output:

\begin{itemize}
    \item Type 1/2: hard stool pattern
    \item Type 3/4: normal-range stool pattern
    \item Type 5/6/7: loose-to-watery stool pattern
\end{itemize}

At the same time, Type 7 probability should be preserved as an independent watery-risk signal. This keeps the product output simple while retaining clinically important risk information.

\section{Current Best Model}

The current primary model is:

\begin{verbatim}
checkpoints_clean_split_convnext_tiny/bsfs_convnext_tiny_final.pth
\end{verbatim}

It uses:

\begin{itemize}
    \item Backbone: ConvNeXt-Tiny
    \item Input size: 300 x 300
    \item Evaluation: deterministic 5-view TTA
    \item Dataset: GuTelligence-StoMy-Clean-Split
\end{itemize}

\section{Current Best Performance}

\subsection{Single Clean Test Set}

\begin{table}[h]
\centering
\begin{tabular}{lr}
\toprule
Metric & Value \\
\midrule
Strict 7-class accuracy & 48.15\% \\
Type 3/4 relaxed accuracy & 57.41\% \\
4-class grouped accuracy & 70.37\% \\
3-class grouped accuracy & 79.63\% \\
Macro F1, strict 7-class & 0.4727 \\
Adjacent accuracy & 86.57\% \\
MAE & 0.7130 \\
\bottomrule
\end{tabular}
\caption{Best single clean-test performance from the current ConvNeXt-Tiny model.}
\end{table}

\subsection{Repeated Clean Split Estimate}

Repeated split evaluation gives a more conservative estimate of generalization:

\begin{table}[h]
\centering
\begin{tabular}{lrr}
\toprule
Metric & Mean & Std \\
\midrule
Strict 7-class accuracy & 46.45\% & 7.77\% \\
4-class grouped accuracy & 64.20\% & 5.86\% \\
3-class grouped accuracy & 72.38\% & 2.55\% \\
Macro F1, strict 7-class & 0.4188 & 0.0823 \\
Adjacent accuracy & 79.17\% & 3.79\% \\
\bottomrule
\end{tabular}
\caption{Repeated clean split results across three leakage-safe splits.}
\end{table}

\section{Completed Experiments and Decisions}

\subsection{Data Leakage Control}

The original dataset split was found to have duplicate or near-duplicate leakage risk. A clean split was created, and all serious model comparisons now use this clean split or repeated clean split manifests.

\subsection{Architecture Experiments}

The following model directions were tested and rejected as replacements:

\begin{itemize}
    \item EfficientNet-B4 tuned baselines
    \item High-resolution EfficientNet-B4
    \item Direct 4-class ConvNeXt
    \item Direct 6-class Type 3/4 merged ConvNeXt
    \item EfficientNetV2-S
    \item ConvNeXt-Small
\end{itemize}

The key finding is that larger or stronger backbones did not improve repeated-split generalization. ConvNeXt-Tiny remains the best overall candidate.

\subsection{Specialist and Multi-Task Experiments}

The following approaches were also tested and rejected:

\begin{itemize}
    \item Type 7 binary gate
    \item Type 7 cascade
    \item Type 2/3/4 neighbor classifier
    \item Ordinal soft-label fine-tuning
    \item Ordinal-aware ConvNeXt loss
    \item 4-class group plus 7-class raw multi-task learning
    \item Raw-primary multi-task learning with auxiliary group loss
\end{itemize}

These experiments did not reliably improve grouped accuracy, strict accuracy, macro F1, or Type 7 recall across repeated splits.

\section{Interpretation}

The main blocker is unlikely to be model capacity alone. The evidence points toward:

\begin{itemize}
    \item label ambiguity between adjacent BSFS classes;
    \item limited Type 7 support;
    \item public/Roboflow-style images not matching the future smart-seat camera distribution;
    \item weak generalization across leakage-safe repeated splits;
    \item mismatch between strict 7-class accuracy and the product's clinically meaningful output.
\end{itemize}

For product use, exact 7-class prediction is probably too brittle at this stage. A grouped and ordinal interpretation is more aligned with both the BSFS scale and the smart toilet use case.

\section{Recommended Product Output}

The recommended near-term output format is:

\begin{verbatim}
primary_group: hard | normal_range | loose_watery
primary_group_confidence: float
bsfs_continuous_score: float
type7_probability: float
watery_risk_flag: low | medium | high
\end{verbatim}

The primary output should be 3-class grouped:

\begin{table}[h]
\centering
\begin{tabular}{ll}
\toprule
Product Group & BSFS Types \\
\midrule
Hard & Type 1/2 \\
Normal range & Type 3/4 \\
Loose-to-watery & Type 5/6/7 \\
\bottomrule
\end{tabular}
\caption{Recommended product-facing BSFS grouping.}
\end{table}

Type 7 should not disappear. It should remain as a separate probability and risk signal, even when the main group is loose-to-watery.

\section{Next Research Direction}

\subsection{Phase A: Formal 3-Class Evaluation}

Create a formal evaluation script for:

\begin{itemize}
    \item 3-class grouped accuracy;
    \item 3-class macro F1;
    \item confidence-bin accuracy;
    \item Type 7 probability metrics;
    \item Type 7 AUROC and recall at selected thresholds;
    \item watery-risk false negative rate.
\end{itemize}

Baseline target:

\begin{itemize}
    \item Current single clean-test 3-class accuracy: 79.63\%
    \item Current repeated-split 3-class accuracy mean: 72.38\%
\end{itemize}

\subsection{Phase B: Product Confidence and Calibration}

The model currently produces softmax probabilities, but these probabilities are not yet calibrated. The next step is to measure and improve:

\begin{itemize}
    \item reliability curves;
    \item expected calibration error;
    \item confidence threshold behavior;
    \item low-confidence abstention;
    \item Type 7 risk thresholding.
\end{itemize}

\subsection{Phase C: 3-Class Primary Model with Type 7 Risk Head}

After formal post-hoc evaluation, train a model with:

\begin{itemize}
    \item a 3-class primary group head;
    \item a Type 7 binary risk head;
    \item optional raw 7-class auxiliary head for diagnostic transparency.
\end{itemize}

This architecture better matches the product objective:

\begin{verbatim}
main task: hard vs normal vs loose_watery
risk task: Type 7 watery probability
\end{verbatim}

\subsection{Phase D: Device-Native Validation}

To move toward a 90\% product-level target, the project needs a validation set collected from the actual smart toilet seat camera module. This should include:

\begin{itemize}
    \item native camera angle;
    \item native lighting;
    \item toilet bowl background;
    \item real capture timing;
    \item multiple frames per event if available.
\end{itemize}

Without device-native data, model-only gains on the current public-style dataset are unlikely to reach product-grade reliability.

\section{Proposed Success Criteria}

\begin{table}[h]
\centering
\begin{tabular}{ll}
\toprule
Metric & Near-Term Target \\
\midrule
Single clean-test 3-class accuracy & $\geq$ 82\% \\
Repeated-split 3-class accuracy mean & $\geq$ 75\% \\
Type 7 AUROC & to be established \\
Type 7 recall at risk threshold & to be established \\
Watery-risk false negative rate & minimize and report explicitly \\
Device-native validation accuracy & to be established \\
\bottomrule
\end{tabular}
\caption{Recommended next-stage success criteria.}
\end{table}

\section{Conclusion}

The current best product-aligned result is 79.63\% 3-class grouped accuracy on the clean test set, with a more conservative repeated-split mean of 72.38\%. This is the most promising direction so far.

The recommended next step is not another larger backbone. Instead, the project should formalize the 3-class product output, preserve Type 7 probability as a separate risk signal, calibrate confidence, and validate on device-native data.

\end{document}
